\section{实直线上的连通子空间}\label{section:ch03-connected-subspaces-real}

上一节的定理给出了用已知的连通空间来构造新的连通空间的方法。那么，从哪里开始去寻找一些连通空间呢？其中最理想的莫过于实直线了。我们将证明 $\mathbb{R}$ 是连通的，以及 $\mathbb{R}$ 中的每一个区间和射线都是连通的。

一个应用就是适当地推广微积分中的介值定理。另一个应用就是证明欧氏空间中诸如球和球面那些熟知的空间的连通性。这个证明涉及我们将要讨论的道路连通性这样一个新的概念。

$\mathbb{R}$ 中的区间和射线的连通性大家可能从数学分析中早已很熟悉了。这里，我们在更广泛的意义下给出证明，并指明这件事与 $\mathbb{R}$ 的代数性质无关，而仅与它的序的性质有关。为了明确起见，我们将对任意与 $\mathbb{R}$ 有同样的序性质的那些全序集证明这个定理。这种集合称为线性连续统。

\begin{definition}\label{definition:ch03-linear-continuum}
若 $L$ 是多于一个元素的全序集，并且满足条件：
\begin{enumerate}
\item $L$ 有上确界性质.
\item 若 $x<y$，则存在 $z$ 使得 $x<z<y$.
\end{enumerate}
则称 $L$ 是一个\textbf{线性连续统}（linear continuum）。
\end{definition}

\begin{theorem}\label{theorem:ch03-linear-continuum-connected}
若 $L$ 是一个赋予序拓扑的线性连续统，则 $L$ 是连通的，并且 $L$ 的每一个区间和每一条射线也都是连通的。
\end{theorem}

\begin{proof}
曾经给出过以下定义：称 $L$ 的一个子空间 $Y$ 是凸的，是指对于满足 $a<b$ 的 $Y$ 中的任何一个点对 $a$ 和 $b$，$L$ 中的点组成的区间 $[a,b]$ 包含于 $Y$。以下我们证明：若 $Y$ 是 $L$ 的一个凸子集，则 $Y$ 是连通的。

设 $Y$ 是两个非空无交集合 $A$ 和 $B$ 之并，其中 $A$ 和 $B$ 都是 $Y$ 的开集。选取 $a\in A$ 和 $b\in B$，并且为讨论方便，假设 $a<b$。于是 $L$ 中的区间 $[a,b]$ 包含于 $Y$。于是，$[a,b]$ 可以表示成无交集合
\begin{align*}
A_{0}=A\cap[a,b]\quad\text{和}\quad B_{0}=B\cap[a,b]
\end{align*}
之并，并且每一集合就子空间的序拓扑而言都是 $[a,b]$ 的开集。由 $a\in A_{0}$ 及 $b\in B_{0}$ 知它们都是非空的。因此 $A_{0}$ 和 $B_{0}$ 是 $[a,b]$ 的一个分割。

令 $c=\sup A_{0}$。以下证明 $c$ 既不属于 $A_{0}$ 也不属于 $B_{0}$。这与 $[a,b]$ 等于 $A_{0}$ 与 $B_{0}$ 的并矛盾。

情形 1. 假定 $c\in B_{0}$。则 $c\neq a$，于是或者 $c=b$，或者 $a<c<b$。无论哪种情况，由于 $B_{0}$ 是 $[a,b]$ 中的开集，总存在一个形如 $(d,c]$ 的区间包含于 $B_{0}$。若 $c=b$，立刻得出矛盾，因为这时 $d$ 是 $A_{0}$ 的一个小于 $c$ 的上界。如果 $c<b$，我们有 $(c,b]$ 与 $A_{0}$ 无交（由于 $c$ 是 $A_{0}$ 的一个上界）。于是
\begin{align*}
(d,b]=(d,c]\cup(c,b]
\end{align*}
与 $A_{0}$ 无交。由此可见 $d$ 仍然是 $A_{0}$ 的一个小于 $c$ 的上界，与假设矛盾。参见\cref{figure:ch03-linear-case1}。

情形 2. 假定 $c\in A_{0}$。则 $c\neq b$，于是或者 $c=a$，或者 $a<c<b$。因为 $A_{0}$ 在 $[a,b]$ 中是开的，一定存在某一个形如 $[c,e)$ 的区间包含在 $A_{0}$ 中。参见\cref{figure:ch03-linear-case2}。根据线性连续统 $L$ 的序性质 (2)，可在 $L$ 中选取一点 $z$，使得 $c<z<e$。于是 $z\in A_{0}$，这与 $c$ 是 $A_{0}$ 的一个上界矛盾。
\end{proof}

\par
\begin{figure}[htbp]
\centering
\includegraphics[width=0.45\textwidth]{ch03-fig24-1.png}
\caption{情形 1}\label{figure:ch03-linear-case1}
\end{figure}

\par
\begin{figure}[htbp]
\centering
\includegraphics[width=0.45\textwidth]{ch03-fig24-2.png}
\caption{情形 2}\label{figure:ch03-linear-case2}
\end{figure}

\begin{corollary}\label{corollary:ch03-real-line-connected}
实直线 $\mathbb{R}$ 是连通的，并且 $\mathbb{R}$ 中的每一个区间和射线也都是连通的。
\end{corollary}

作为一个应用，我们证明微积分中的介值定理的一个适当的推广形式。

\begin{theorem}[介值定理]\label{theorem:ch03-intermediate-value}
设 $f\colon X\to Y$ 是从连通空间 $X$ 到具有序拓扑的全序集 $Y$ 的一个连续映射。若 $a$ 和 $b$ 是 $X$ 的两个点并且 $r$ 是 $Y$ 中介于 $f(a)$ 和 $f(b)$ 之间的一个点\footnote{根据定理的证明可见，``$r$ 是 $Y$ 中介于 $f(a)$ 和 $f(b)$ 之间的一个点''意思是指或者 $f(a)<r<f(b)$ 成立，或者 $f(b)<r<f(a)$ 成立。}，则 $X$ 中存在一个点 $c$ 使得 $f(c)=r$。
\end{theorem}

作为特殊情形，若取 $X$ 为 $\mathbb{R}$ 中的闭区间并且取 $Y$ 为 $\mathbb{R}$，便可得到微积分中的介值定理。

\begin{proof}
假定定理所设条件成立。那么
\begin{align*}
A=f(X)\cap(-\infty,r)\quad\text{和}\quad B=f(X)\cap(r,+\infty)
\end{align*}
是无交的，由于 $A$ 和 $B$ 中一个包含 $f(a)$，另一个包含 $f(b)$，所以它们都是非空的。由于它们是 $Y$ 的开射线与 $f(X)$ 的交，所以都是 $f(X)$ 的开集。若 $X$ 中不存在使得 $f(c)=r$ 的点 $c$，那么 $f(X)$ 就等于是集合 $A$ 与 $B$ 的并。于是 $A$ 和 $B$ 便组成 $f(X)$ 的一个分割，这与连通空间的连续像是连通的这一事实相矛盾。
\end{proof}

\begin{example}\label{example:ch03-ordered-rectangle}
一个异于 $\mathbb{R}$ 的线性连续统的例子是有序矩形。我们只验证上确界性质（线性连续统定义中的第二个要求是显然的）。设 $A$ 为 $I\times I$ 的一个子集。设 $\pi_{1}\colon I\times I\to I$ 是向第一个坐标空间的投射。并令 $b=\sup\pi_{1}(A)$。若 $b\in\pi_{1}(A)$，则 $A$ 与 $I\times I$ 的子集 $b\times I$ 相交，因为 $b\times I$ 具有 $I$ 的序型，$A\cap(b\times I)$ 有上确界 $b\times c$，它也是 $A$ 的上确界。参见\cref{figure:ch03-ordered-rectangle-fig}。如果 $b\notin\pi_{1}(A)$，则 $b\times 0$ 是 $A$ 的上确界。而满足 $b'<b$ 的形如 $b'\times c$ 的元素都不可能是 $A$ 的上界，因为此时 $b'$ 将是 $\pi_{1}(A)$ 的一个上界。
\end{example}

\par
\begin{figure}[htbp]
\centering
\includegraphics[width=0.6\textwidth]{ch03-fig24-3.png}
\caption{有序矩形的上确界性质}\label{figure:ch03-ordered-rectangle-fig}
\end{figure}

\begin{example}\label{example:ch03-well-ordered-product}
如果 $X$ 是一个良序集，留给读者自行验证 $X\times[0,1)$ 关于字典序是一个线性连续统。这个集合可以看成是在 $X$ 的每一个元素和它的紧接后元之间``填充''一个具有 $(0,1)$ 序型的集合而构成的。
\end{example}

$\mathbb{R}$ 中的区间的连通性引出了一个判定空间 $X$ 连通的特别有用的准则，即 $X$ 中的任何一对点都能用 $X$ 中的一条道路连接。

\begin{definition}\label{definition:ch03-path}
设 $x$ 与 $y$ 为空间 $X$ 的两点，$X$ 中从 $x$ 到 $y$ 的一条\textbf{道路}（path）是指从实直线的某一个闭区间 $[a,b]$ 到 $X$ 的一个连续映射 $f\colon[a,b]\to X$，使得 $f(a)=x$ 和 $f(b)=y$。如果空间 $X$ 中每一对点都能用 $X$ 中的一条道路连接，则称 $X$ 是\textbf{道路连通的}（path connected）。
\end{definition}

\begin{proposition}\label{proposition:ch03-path-connected-implies-connected}
道路连通空间 $X$ 必然连通。
\end{proposition}

\begin{proof}
设 $X=A\cup B$ 是 $X$ 的一个分割。令 $f\colon[a,b]\to X$ 为 $X$ 中任何一条道路。而集合 $f([a,b])$ 是连通集的连续像，所以它是连通的。于是它必定完全包含于 $A$ 或者包含于 $B$ 中。因此，$X$ 中便不存在连接 $A$ 中的点与 $B$ 中的点的道路，这与 $X$ 是道路连通的假设相矛盾。
\end{proof}

上述命题的逆命题并不成立，一个连通空间未必是道路连通的。见下面的例\,\cref{example:ch03-ordered-rectangle-not-path} 和例\,\cref{example:ch03-topologist-sine}。

\begin{example}\label{example:ch03-unit-ball}
$\mathbb{R}^{n}$ 中的\textbf{单位球}（unit ball）$B^{n}$ 定义为
\begin{align*}
B^{n}=\{\,\boldsymbol{x}\mid\|\boldsymbol{x}\|\leqslant 1\,\},
\end{align*}
其中
\begin{align*}
\|\boldsymbol{x}\|=\|(x_{1},\cdots,x_{n})\|=(x_{1}^{2}+\cdots+x_{n}^{2})^{1/2}.
\end{align*}
单位球是道路连通的。事实上，任意给定 $B^{n}$ 中两点 $\boldsymbol{x}$ 与 $\boldsymbol{y}$，由
\begin{align*}
f(t)=(1-t)\boldsymbol{x}+t\boldsymbol{y}
\end{align*}
定义的直线路 $f\colon[0,1]\to\mathbb{R}^{n}$ 完全含在 $B^{n}$ 中。因为若 $\boldsymbol{x}$ 与 $\boldsymbol{y}$ 在 $B^{n}$ 中并且 $t$ 在 $[0,1]$ 中，那么
\begin{align*}
\|f(t)\|\leqslant(1-t)\|\boldsymbol{x}\|+t\|\boldsymbol{y}\|\leqslant 1.
\end{align*}
类似地，可以证明 $\mathbb{R}^{n}$ 中每一个开球 $B_{d}(\boldsymbol{x},\varepsilon)$ 和每一个闭球 $\overline{B}_{d}(\boldsymbol{x},\varepsilon)$ 都是道路连通的。
\end{example}

\begin{example}\label{example:ch03-punctured}
\textbf{穿孔欧氏空间}（punctured Euclidean space）定义为空间 $\mathbb{R}^{n}\setminus\{\boldsymbol{0}\}$，其中 $\boldsymbol{0}$ 是 $\mathbb{R}^{n}$ 中的原点。对于 $n>1$，它是道路连通的：任意给定异于 $\boldsymbol{0}$ 的两点 $\boldsymbol{x}$ 与 $\boldsymbol{y}$，如果它们的连线不通过原点，就用直线路连接 $\boldsymbol{x}$ 与 $\boldsymbol{y}$。否则，我们可选取一点 $\boldsymbol{z}$ 不在 $\boldsymbol{x}$ 与 $\boldsymbol{y}$ 连线上，连一条从 $\boldsymbol{x}$ 到 $\boldsymbol{z}$，再从 $\boldsymbol{z}$ 到 $\boldsymbol{y}$ 的折线路。
\end{example}

\begin{example}\label{example:ch03-unit-sphere}
$\mathbb{R}^{n}$ 中的\textbf{单位球面}（unit sphere）$S^{n-1}$ 定义为
\begin{align*}
S^{n-1}=\{\,\boldsymbol{x}\mid\|\boldsymbol{x}\|=1\,\}.
\end{align*}
当 $n>1$ 时它是道路连通的。因为用 $g(\boldsymbol{x})=\boldsymbol{x}/\|\boldsymbol{x}\|$ 所定义的映射 $g\colon\mathbb{R}^{n}\setminus\{\boldsymbol{0}\}\to S^{n-1}$ 是连续满射。易见，一个道路连通空间的连续像是道路连通的。
\end{example}

\begin{example}\label{example:ch03-ordered-rectangle-not-path}
有序矩形 $I_{o}^{2}$ 是连通的，但不是道路连通的。

因为 $I_{o}^{2}$ 是一个线性连续统，所以它是连通的。设 $p=0\times 0$，$q=1\times 1$，如果存在一个连接 $p$ 和 $q$ 的道路 $f\colon[a,b]\to I_{o}^{2}$，必定导致矛盾。根据介值定理，像集 $f([a,b])$ 必定包含 $I_{o}^{2}$ 的每一个点 $x\times y$。因此，对于每一个点 $x\in I$，集合
\begin{align*}
U_{x}=f^{-1}(x\times(0,1))
\end{align*}
是 $[a,b]$ 中一个非空子集。并且根据连续性，它是 $[a,b]$ 中的一个开集。参见\cref{figure:ch03-ordered-rectangle-not-path-fig}。对于每一个 $x\in I$，在 $U_{x}$ 中选取一个有理数 $q_{x}$。因为这些集合 $U_{x}$ 无交，映射 $x\mapsto q_{x}$ 是一个从 $I$ 到 $\mathbb{Q}$ 中的单射。这与区间 $I$ 不可数矛盾（后面我们将证明 $I$ 是不可数的）。
\end{example}

\par
\begin{figure}[htbp]
\centering
\includegraphics[width=0.5\textwidth]{ch03-fig24-4.png}
\caption{有序矩形不是道路连通的}\label{figure:ch03-ordered-rectangle-not-path-fig}
\end{figure}

\begin{example}\label{example:ch03-topologist-sine}
设 $S$ 表示平面上的下列子集：
\begin{align*}
S=\{\,x\times\sin(1/x)\mid 0<x\leqslant 1\,\}.
\end{align*}
因为 $S$ 是连通集 $(0,1]$ 的连续像，所以 $S$ 是连通的。因此它在 $\mathbb{R}^{2}$ 中的闭包 $\overline{S}$ 也是连通的。集合 $\overline{S}$ 是拓扑学上的一个经典例子，称为\textbf{拓扑学家的正弦曲线}（topologist's sine curve）。如\cref{figure:ch03-sine-fig}所示。它是 $S$ 与一个垂直区间 $0\times[-1,1]$ 之并。我们来证明 $\overline{S}$ 不是道路连通的。

假设 $f\colon[a,c]\to\overline{S}$ 是一个连接原点与 $S$ 中一点的道路。则所有满足 $f(t)\in 0\times[-1,1]$ 的 $t$ 构成一个闭集，从而有最大元 $b$。那么 $f\colon[b,c]\to\overline{S}$ 是一个将 $b$ 映到 $0\times[-1,1]$ 中，将 $[b,c]$ 中所有异于 $b$ 的点映到 $S$ 中的道路。

为讨论方便，以 $[0,1]$ 代替 $[b,c]$，并且记 $f(t)=(x(t),y(t))$。则 $x(0)=0$，且当 $t>0$ 时，$x(t)>0$，$y(t)=\sin(1/x(t))$。我们来证明存在点的一个序列 $t_{n}\to 0$ 使得 $y(t_{n})=(-1)^{n}$。由于序列 $y(t_{n})$ 不收敛，从而与 $f$ 的连续性矛盾。

我们可按照以下方式选取 $t_{n}$：对于给定的 $n$，选取满足 $0<u<x(1/n)$ 的 $u$ 使得 $\sin(1/u)=(-1)^{n}$。那么由介值定理知，存在满足 $0<t_{n}<1/n$ 的 $t_{n}$ 使得 $x(t_{n})=u$。
\end{example}

\par
\begin{figure}[htbp]
\centering
\includegraphics[width=0.5\textwidth]{ch03-fig24-5.png}
\caption{拓扑学家的正弦曲线}\label{figure:ch03-sine-fig}
\end{figure}

\subsection*{习题}

\begin{enumerate}
\item (a) 证明：空间 $(0,1)$，$(0,1]$ 与 $[0,1]$ 彼此不同胚。[提示：从每一空间中去掉一个点将发生什么情况？]
\item (b) 设 $f\colon X\to Y$ 和 $g\colon Y\to X$ 都是嵌入映射。举例说明 $X$ 和 $Y$ 未必同胚。
\item (c) 证明：若 $n>1$，则 $\mathbb{R}^{n}$ 与 $\mathbb{R}$ 不同胚。
\item 设 $f\colon S^{1}\to\mathbb{R}$ 是一个连续映射。证明存在 $S^{1}$ 的一个点 $x$，使得 $f(x)=f(-x)$。
\item 设 $f\colon X\to X$ 是连续的。证明：若 $X=[0,1]$，则存在一点 $x$ 使得 $f(x)=x$。点 $x$ 称之为 $f$ 的一个\textbf{不动点}（fixed point）。若 $X$ 为 $[0,1)$ 或 $(0,1)$，结论如何？
\item 设 $X$ 是赋予序拓扑的一个全序集。证明：若 $X$ 连通，则 $X$ 为线性连续统。
\item 讨论下列具有字典序的集合，其中哪些是线性连续统？
\begin{enumerate}
\item $\mathbb{Z}_{+}\times[0,1)$.
\item $[0,1)\times\mathbb{Z}_{+}$.
\item $[0,1)\times[0,1]$.
\item $[0,1]\times[0,1)$.
\end{enumerate}
\item 证明：若 $X$ 是一个良序集，则 $X\times[0,1)$ 关于字典序是一个线性连续统。
\item (a) 设 $X$ 与 $Y$ 都是赋予序拓扑的全序集。证明：若 $f\colon X\to Y$ 是一个保序满射，则 $f$ 为同胚。
\item (b) 设 $X=Y=\overline{\mathbb{R}}_{+}$。对于取定的正整数 $n$，证明：函数 $f(x)=x^{n}$ 是一个保序满射。因而其逆映射，即求 $n$ 次方根，是连续的。
\item (c) 设 $X$ 是 $\mathbb{R}$ 的一个子空间 $(-\infty,-1)\cup[0,\infty)$。函数 $f\colon X\to\mathbb{R}$ 定义为：当 $x<-1$ 时，$f(x)=x+1$。当 $x\geqslant 0$ 时，$f(x)=x$。证明：$f$ 是一个保序满射。$f$ 是否是一个同胚？与 (a) 小题中的结论进行比较。
\item (a) 道路连通空间的积空间必定道路连通吗？
\item (b) 若 $A\subset X$ 且 $A$ 是道路连通的，$\overline{A}$ 必定道路连通吗？
\item (c) 若 $f\colon X\to Y$ 连续，并且 $X$ 道路连通，$f(X)$ 必定道路连通吗？
\item (d) 若 $\{A_{\alpha}\}$ 是 $X$ 的道路连通子空间的族，并且 $\bigcap A_{\alpha}\neq\varnothing$，$\bigcup A_{\alpha}$ 必定道路连通吗？
\item 设 $\mathbb{R}$ 不可数。证明：若 $A$ 为 $\mathbb{R}^{2}$ 的可数字集，那么 $\mathbb{R}^{2}\setminus A$ 是道路连通的。[提示：通过 $\mathbb{R}^{2}$ 中一个定点的直线有多少条？]
\item 证明：若 $U$ 是 $\mathbb{R}^{2}$ 的一个连通开子空间，则 $U$ 道路连通。[提示：对于给定的 $x_{0}\in U$，证明在 $U$ 中可以与 $x_{0}$ 道路连接的那些点的集合是 $U$ 中既开又闭的。]
\item 若 $A$ 是 $X$ 的一个连通子空间，$\mathrm{Int}\,A$ 和 $\mathrm{Bd}\,A$ 必定连通吗？其逆成立吗？验证你的结论。
\item 如前，$S_{\Omega}$ 表示极小不可数良序集。用 $L$ 表示在具有字典序的全序集 $S_{\Omega}\times[0,1)$ 中去掉它的最小元后所得到的集合。集合 $L$ 是拓扑学中的一个经典例子，称之为\textbf{长线}（long line）。

\begin{theorem}
长线是道路连通的，并且局部同胚于 $\mathbb{R}$，但它不能嵌入到 $\mathbb{R}$ 中。
\end{theorem}
\begin{enumerate}
\item (a) 设 $X$ 是一个全序集，并且 $a<b<c$ 都是 $X$ 中的点。证明：$[a,c)$ 具有 $[0,1)$ 的序型当且仅当 $[a,b)$ 和 $[b,c)$ 都有 $[0,1)$ 的序型。
\item (b) 设 $X$ 是一个全序集，$x_{0}<x_{1}<\cdots$ 是 $X$ 中的一个递增序列。设 $b=\sup\{x_{i}\}$。证明：$[x_{0},b)$ 具有 $[0,1)$ 的序型当且仅当每一个区间 $[x_{i},x_{i+1})$ 都具有 $[0,1)$ 的序型。
\item (c) 设 $a_{0}$ 为 $S_{\Omega}$ 的最小元。对于 $S_{\Omega}$ 的每一异于 $a_{0}$ 的点 $a$，证明 $S_{\Omega}\times[0,1)$ 的每一区间 $[a_{0}\times 0,a\times 0)$ 具有 $[0,1)$ 的序型。[提示：用超限归纳法。或者 $a$ 在 $S_{\Omega}$ 中有紧接前元，或者存在一个 $S_{\Omega}$ 中的严格递增序列 $a_{i}$，使得 $a=\sup\{a_{i}\}$。]
\item (d) 证明 $L$ 是道路连通的。
\item (e) 证明 $L$ 的每一点都存在一个同胚于 $\mathbb{R}$ 中的开区间的邻域。
\item (f) 证明 $L$ 不能嵌入 $\mathbb{R}$ 中，实际上对于任何 $n$，它也不能嵌入 $\mathbb{R}^{n}$ 中。[提示：$\mathbb{R}^{n}$ 的任何子空间都有一个相对于子空间拓扑的可数基。]
\end{enumerate}
\end{enumerate}